% Generated draft: requires mathematical and source review.
\documentclass[11pt]{article}
\usepackage[margin=25mm]{geometry}
\usepackage{fontspec}
\setmainfont{Latin Modern Roman}
\usepackage{amsmath,amssymb,mathtools,unicode-math}
\setmathfont{Latin Modern Math}
\usepackage{enumitem,needspace,xurl,hyperref}
\setlength{\parindent}{0pt}
\setlength{\parskip}{5pt}
\setlength{\emergencystretch}{3em}
\newcommand{\N}{\mathbb{N}}
\newcommand{\Z}{\mathbb{Z}}
\newcommand{\Q}{\mathbb{Q}}
\newcommand{\R}{\mathbb{R}}
\newcommand{\C}{\mathbb{C}}
\newcommand{\F}{\mathbb{F}}
\newcommand{\A}{\mathbb{A}}
\newcommand{\PP}{\mathbb{P}}
\newcommand{\E}{\mathbb{E}}
\newcommand{\eps}{\varepsilon}
\newcommand{\dd}{\,\mathrm{d}}
\begin{document}
\section*{UnsolvedMath ID: 800855}
\textbf{Source problem code:} catalogue entry
\subsection{Definitions and mathematical statement}
For a prime \(p\equiv 1\pmod 4\), the Paley graph \(P_p\) has vertex set \(\F_p\); distinct vertices \(x,y\) are adjacent precisely when \(x-y\) is a nonzero square in \(\F_p\). For a finite simple graph \(H\), its complement \(\overline H\) has the same vertex set and has an edge between distinct vertices precisely when \(H\) does not. For a nonempty finite simple graph \(H\), its Lov\'asz theta number is
\[
\vartheta(H)=\max\left\{\sum_{u,v\in V(H)}X_{uv}: X\in\R^{V(H)\times V(H)},\ X=X^{\mathsf T},\ X\succeq0,\ \operatorname{tr}(X)=1,\ X_{uv}=0\ \text{for every }\{u,v\}\in E(H)\right\}.
\]
For the graph with no vertices, set \(\vartheta(H)=0\).

The supplied record denotes by \(G_{p,2}\) the ``2-localization'' of the Paley graph and refers to a description ``above.'' That description is absent from the supplied record. Consequently, the vertex set, adjacency relation, and any choices or additional quantifiers involved in \(G_{p,2}\) remain unspecified; the definition of \(P_p\) alone does not resolve them.

\textbf{Conjecture 3, conditional on the missing definition.} With \(G_{p,2}\) interpreted exactly as in the omitted source description, the proposed assertion is that there exists \(p_0\in\N\) such that, for every prime \(p\ge p_0\) satisfying \(p\equiv1\pmod4\),
\[
\vartheta\!\left(\overline{G_{p,2}}\right)\le \frac{2}{3}\sqrt p.
\]
This is an incomplete formulation until the source's definition of 2-localization, including any choices and their quantification, is recovered. The registry labels the problem open; that label is supplied metadata, not an independently verified mathematical or literature assessment.
\subsection{Short English statement}
The proposed bound says that, for every sufficiently large prime \(p\equiv1\pmod4\), the complement of the specified 2-localization of the Paley graph has Lov\'asz theta number at most \(\frac23\sqrt p\). The supplied record omits the definition of 2-localization, so the graph and any quantification over its defining choices require clarification.
\subsection{Sources}
\begin{itemize}
\item Supplied local UnsolvedMath registry record, permanent ID \(800855\), exact source problem code \texttt{catalogue entry}; upstream ID \(800855\), revision \texttt{d371ee917ac9647fbb06cd6fc529deec6f3dda82}. The record attributes the proposal to Bandeira et al., gives proposed year 2026, and labels its status open.
\item Bandeira et al., \emph{Randomstrasse101: Open Problems of 2025} (2026), Conjecture 3. Supplied source URL: \url{https://arxiv.org/abs/2603.29571}. The record reports access on 2026-10-16 and extraction by \texttt{arxiv\_latex}. These are supplied provenance details; neither the source nor its access history has been independently verified here. The supplied excerpt refers to an earlier definition of 2-localization that is not included in the record.
\end{itemize}
\end{document}
